% Source: https://artofproblemsolving.com/community/c3h1377788_difference_of_consecutive_squares
% Stable identity: aops-topic-1377788-q01
Factor: $a^2-b^2=(a-b)(a+b)$.

(a) Take $a=3$, $b=1$: $a^2-b^2=9-1=8$, but $a+b=4$. So the claim is false.

(b) Since $a+b>0$, the equation $(a-b)(a+b)=a+b$ is equivalent to $a-b=1$. So the equation holds \emph{exactly} when $a$ and $b$ are consecutive integers, $\boxed{a=b+1}$; both of the student's examples are of this kind. In general $a^2-b^2$ is $(a-b)$ times $a+b$.
